% Green V2X — Full chat transcript as project overview (verbatim content, LaTeX presentation)
\documentclass[11pt,a4paper]{article}

\usepackage[utf8]{inputenc}
\usepackage[T1]{fontenc}
\usepackage{lmodern}
\usepackage[margin=2.2cm]{geometry}
\usepackage{xcolor}
\usepackage{hyperref}
\usepackage{booktabs}
\usepackage{array}
\usepackage{enumitem}
\usepackage{tcolorbox}
\usepackage{fancyhdr}
\usepackage{titlesec}
\usepackage{microtype}

\definecolor{userbg}{RGB}{232,242,255}
\definecolor{assistbg}{RGB}{245,250,245}
\definecolor{accent}{RGB}{0,90,50}
\definecolor{ruleline}{RGB}{180,200,190}

\hypersetup{colorlinks=true,linkcolor=accent,urlcolor=accent,citecolor=accent}

\tcbuselibrary{skins,breakable}

% --- User message box ---
\newtcolorbox{usermsg}{
  enhanced,
  breakable,
  colback=userbg,
  colframe=accent!60!black,
  boxrule=0.8pt,
  arc=3pt,
  left=10pt,right=10pt,top=8pt,bottom=8pt,
  fonttitle=\bfseries\color{accent},
  title={User},
}

% --- Assistant message box ---
\newtcolorbox{assistmsg}{
  enhanced,
  breakable,
  colback=assistbg,
  colframe=ruleline,
  boxrule=0.6pt,
  arc=3pt,
  left=10pt,right=10pt,top=8pt,bottom=8pt,
  fonttitle=\bfseries\color{accent!80!black},
  title={Assistant},
}

\titleformat{\section}{\Large\bfseries\color{accent}}{}{0em}{}[\titlerule[0.4pt]]
\titleformat{\subsection}{\large\bfseries}{}{0em}{}

\pagestyle{fancy}
\fancyhf{}
\fancyhead[L]{\small\textit{Green V2X — Conversation overview}}
\fancyhead[R]{\small\thepage}
\renewcommand{\headrulewidth}{0.4pt}
\renewcommand{\footrulewidth}{0pt}

\setlist[itemize]{itemsep=2pt,parsep=0pt,topsep=4pt}

\begin{document}

\begin{center}
{\LARGE\bfseries\color{accent} Green V2X: Energy-Aware Handoff}\\[0.35em]
{\large Project-style overview from full chat transcript}\\[0.5em]
{\small\textit{Content preserved exactly as in the conversation; only layout and \LaTeX{} markup added.}}
\end{center}
\vspace{1em}

% ========== TURN 1 ==========
\begin{usermsg}
FIrst checkout the whole project.

Then tell me what is this project about ?

what is the algorithm here ?

How everything works ?

and also what goal/result is achieving this project ?

I need simple but detailed explanation in easy words please
\end{usermsg}

\begin{assistmsg}
Here's a concise picture of the project after reading the main docs and core code (no \texttt{git checkout} was needed—you meant ``look over'' the repo).

\subsection*{What this project is about}

\textbf{Green V2X} is a \textbf{computer simulation} of cars on a \textbf{highway} talking to \textbf{cell towers (base stations)} over the \textbf{uplink} (vehicle $\rightarrow$ tower). In real life, a car must \textbf{choose which tower to use} as it moves; that choice is called \textbf{handoff}.

This repo implements \textbf{several handoff rules}, especially one that tries to \textbf{save radio energy} (and then \textbf{CO\textsubscript{2}}, using grid electricity intensity). It is \textbf{not} a live network stack---it's a \textbf{research-style discrete-time simulator} that steps time, moves vehicles, updates radio, counts handoffs, and totals energy and throughput.

\subsection*{What the main algorithm is (energy-aware handoff)}

The \textbf{star algorithm} is \textbf{\texttt{EnergyAwareHandoff}} in \texttt{src/algorithms/energy\_aware\_handoff.py}. In plain terms it does this:

\begin{enumerate}[label=\arabic*.]
\item \textbf{For each candidate tower}, it estimates:
  \begin{itemize}
  \item \textbf{Transmit power} needed for a reliable link (path loss, weather, shadowing).
  \item \textbf{Data rate} from a Shannon-style formula using \textbf{SINR} (signal vs noise + interference).
  \item \textbf{Energy per bit (EPB)} = how many joules of device power you spend per bit sent, using the \textbf{energy model} (PA, circuits, baseband, etc.).
  \end{itemize}
\item It \textbf{combines EPB with tower load}: a tower with more users gets a \textbf{worse score} (higher ``metric''), so the policy prefers \textbf{less loaded} cells when EPB is similar---\textbf{load balancing}.
\item It applies \textbf{QoS filters} (minimum data rate, minimum SINR) so it doesn't pick ``cheap but useless'' links.
\item \textbf{Cell-edge penalty}: if you're far from the tower (near the edge), the \textbf{metric is worse}, so it prefers \textbf{more central / stable} connections when possible.
\item \textbf{When to actually hand off} (\texttt{should\_handoff}) uses:
  \begin{itemize}
  \item \textbf{Escape} if the current tower is \textbf{overloaded}.
  \item \textbf{Urgent} switch if \textbf{SINR is bad} or \textbf{EPB is too high} (bad link).
  \item Otherwise: \textbf{hysteresis}---only switch if the new tower is \textbf{meaningfully better} in energy terms (not tiny gains).
  \item \textbf{Don't} switch to a clearly \textbf{worse SINR} link.
  \item \textbf{Time-to-trigger (TTT)} and \textbf{cooldowns} so cars don't \textbf{ping-pong} between towers.
  \end{itemize}
\end{enumerate}

\textbf{Baselines} (for comparison): strongest \textbf{RSSI}, \textbf{SINR}-based, \textbf{load-aware RSSI}, and \textbf{nearest tower} geographically. These are classic or simple rules; the energy-aware one is the one that \textbf{explicitly optimizes EPB + load + stability}.

\subsection*{How everything works (end-to-end)}

\begin{enumerate}[label=\arabic*.]
\item \textbf{\texttt{SimulationConfig}} (\texttt{simulations/config.py}) sets the world: area size, number of towers, vehicles, speeds, \textbf{highway lanes}, \textbf{weather profile}, \textbf{carbon intensity} (kg CO\textsubscript{2} per kWh), duration, time step, etc.
\item \textbf{\texttt{V2XSimulator}} (\texttt{simulations/simulator.py}):
  \begin{itemize}
  \item Places \textbf{base stations} on a \textbf{grid}.
  \item Spawns \textbf{vehicles} (highway mode: fixed lane centers, lane speeds, optional lane changes).
  \item Each \textbf{time step}: vehicles move, \textbf{shadowing/SINR/interference} are computed (cached per step for fairness), each algorithm decides \textbf{stay or hand off}.
  \item \textbf{Energy} accumulates from \textbf{transmit power $\times$ time}, \textbf{handoff overhead} (extra energy and delay), \textbf{bits sent}, etc.
  \item \texttt{run\_comparison()} runs \textbf{each algorithm} with a \textbf{fresh reset} of vehicle and BS state so results are fair.
  \end{itemize}
\item \textbf{\texttt{main.py}} runs \textbf{multi-seed} experiments, \textbf{paired t-tests} (e.g.\ RSSI vs energy-aware), \textbf{plots}, \textbf{scaling} (more vehicles), \textbf{energy-parameter sensitivity}, \textbf{scenario diversity} (weather / urban), and optional \textbf{Bangladesh grid} CO\textsubscript{2} numbers.
\end{enumerate}

So: \textbf{config $\rightarrow$ build map $\rightarrow$ simulate time $\rightarrow$ policies differ $\rightarrow$ metrics compared}.

\subsection*{Goal and what result it achieves}

\textbf{Goal:} Show that \textbf{handoff can be chosen to reduce uplink energy} (and \textbf{CO\textsubscript{2}}) compared to usual signals-only rules, \textbf{without} necessarily giving up throughput, and ideally with \textbf{fewer unnecessary handoffs} or \textbf{better stability} in some cases.

\textbf{Typical outputs} (under \texttt{results/}, described in the README):

\begin{itemize}
\item \textbf{Energy saving \%} (energy-aware vs RSSI and others), \textbf{energy per bit}, \textbf{total joules}.
\item \textbf{Handoffs}, \textbf{ping-pong}, \textbf{handoff delay}, \textbf{service availability}, \textbf{outage} (coverage gaps vs SINR).
\item \textbf{CO\textsubscript{2}} estimates from energy $\times$ \textbf{grid carbon intensity}.
\item \textbf{Plots} (bars, CDFs, scaling curves, weather sweeps, etc.).
\end{itemize}

The README also notes a \textbf{sanity check}: e.g.\ \textbf{$>$15\% EPB saving vs RSSI} on at least one seed in the paper baseline scenario---this is a \textbf{project-specific success criterion}, not a universal law.

\subsection*{In one sentence:}

This project \textbf{simulates highway V2X}, compares \textbf{how handoff is decided}, and argues that an \textbf{energy- and load-aware} policy can \textbf{use less transmit energy and associated CO\textsubscript{2}} than \textbf{RSSI/SINR/nearest} baselines, while reporting \textbf{full radio and QoS-style statistics}---all in \textbf{software}, for \textbf{research and reporting} rather than deployment in real cars.
\end{assistmsg}

% ========== TURN 2 ==========
\begin{usermsg}
check all the results.

Tell me what are the most last simulated results of this whole project ?
\end{usermsg}

\begin{assistmsg}
Here is what is in your \texttt{results/} folder and what the \textbf{most recently written} outputs say.

\subsection*{What was simulated last (by file time)}

On \textbf{April 11, 2026} (local machine time), artifacts were updated in this order (newest first):

\begin{center}
\begin{tabular}{@{}p{2.6cm}p{4.2cm}p{7.5cm}@{}}
\toprule
\textbf{Time (approx.)} & \textbf{What ran} & \textbf{Main outputs} \\
\midrule
\textbf{$\sim$03:46} & \texttt{sweep\_connectivity\_energy.py} & \texttt{parameter\_sweep.json} (\textbf{newest file}) \\
\textbf{$\sim$02:27} & \texttt{weather\_sweep.py} & \texttt{weather\_sweep.json} + 3 bar PNGs \\
\textbf{$\sim$02:23} & Comprehensive validation & \texttt{comprehensive\_validation.json}, \texttt{extrapolation\_validation.png} \\
\textbf{$\sim$02:13} & Scenario diversity (\texttt{main.py} block) & \texttt{scenario\_diversity\_results.json}, PNG \\
\textbf{$\sim$02:11} & Energy model sensitivity & \texttt{energy\_model\_sensitivity.json/png}, \texttt{energy\_model\_justification.md} \\
\textbf{$\sim$02:05} & Scaling study & \texttt{scaling\_results.json}, \texttt{scaling\_energy\_saving\_vs\_vehicles.png} \\
\textbf{$\sim$01:40} & Baseline suite + plots (\texttt{main.py}) & \texttt{simulation\_results.json} + bar/energy/network PNGs \\
\bottomrule
\end{tabular}
\end{center}

\texttt{pareto\_frontier.png} is older (\textbf{Apr 10}), likely from a notebook.

So the \textbf{single newest numerical result file} is \textbf{\texttt{results/parameter\_sweep.json}} (hyperparameter / connectivity grid search). The \textbf{newest ``full project'' baseline bundle} from \texttt{main.py} is the \textbf{$\sim$01:40} run, with \textbf{\texttt{simulation\_results.json} storing only the last seed} (see below).

\subsection*{Latest main baseline (\texttt{simulation\_results.json}, seed \textbf{1011})}

\texttt{simulation\_results.json} matches \textbf{\texttt{main.py}'s last iteration}: config shows \textbf{seed 1011}, \textbf{paper baseline} (20 vehicles, 16 BSs, 1200 s, clear weather, 2800 m area).

\textbf{Per-algorithm totals (that saved run):}

\begin{center}
\begin{tabular}{@{}lrrrr@{}}
\toprule
\textbf{Algorithm} & \textbf{Total energy (J)} & \textbf{Avg EPB} & \textbf{Handoffs} & \textbf{CO\textsubscript{2} (kg, direct)} \\
\midrule
\textbf{energy\_aware} & \textbf{40,604.84} & \textbf{5.14$\times$10$^{-8}$} & 2,447 & 0.005640 \\
rssi & 41,165.85 & 5.30$\times$10$^{-8}$ & 2,748 & 0.005717 \\
sinr & 41,177.78 & 5.24$\times$10$^{-8}$ & 2,761 & 0.005719 \\
load\_aware\_rssi & 41,146.58 & 5.32$\times$10$^{-8}$ & 2,755 & 0.005715 \\
naive\_nearest & 41,088.76 & 6.45$\times$10$^{-8}$ & 407 & 0.005707 \\
\bottomrule
\end{tabular}
\end{center}

\textbf{Energy-aware vs RSSI on this saved file:}

\begin{itemize}
\item \textbf{$\sim$3.0\%} lower average \textbf{energy per bit (EPB)}
\item \textbf{$\sim$1.36\%} lower \textbf{total energy} and \textbf{CO\textsubscript{2}} (same carbon intensity)
\item \textbf{$\sim$11\%} fewer \textbf{handoffs} than RSSI
\end{itemize}

\textbf{Note:} \texttt{main.py} prints \textbf{multi-seed} statistics in the console; the JSON \textbf{does not} aggregate all five seeds---it reflects the \textbf{last seed} only (here \textbf{1011}), as the README states.

\subsection*{Other latest JSON summaries}

\textbf{Scaling} (\texttt{scaling\_results.json}, denser \texttt{scaling\_scenario}, BS radius \textbf{500 m}):

Energy-aware vs RSSI (EPB saving) \textbf{mean $\pm$ std} over seeds 42, 123, 456:

\begin{itemize}
\item 20 vehicles: \textbf{$\sim$4.93\% $\pm$ 0.14\%}
\item 50: \textbf{$\sim$3.45\% $\pm$ 0.29\%}
\item 100: \textbf{$\sim$3.46\% $\pm$ 0.16\%}
\item 200: \textbf{$\sim$3.28\% $\pm$ 0.16\%}
\end{itemize}

\textbf{Energy model sensitivity} (\texttt{energy\_model\_sensitivity.json}):

EPB saving vs RSSI stays \textbf{$\sim$3.3\%} in the PA efficiency and TX-circuit sweeps (small variation across those knobs).

\textbf{Scenario diversity} (\texttt{scenario\_diversity\_results.json}):

EPB saving vs RSSI (mean over 3 seeds):

\begin{itemize}
\item Clear weather: \textbf{$\sim$1.14\%}
\item Heavy rain: \textbf{$\sim$52\%} (large gap---different channel stress vs RSSI in that profile)
\item Urban dense: \textbf{$\sim$17.2\%}
\end{itemize}

\textbf{Weather sweep} (\texttt{weather\_sweep.json}):

Example for \textbf{clear}: EPB saving mean \textbf{$\sim$1.32\%}, handoff reduction mean \textbf{$\sim$17.3\%}, CO\textsubscript{2} saving mean \textbf{$\sim$0.84\%} (with 95\% CIs in the file for each weather).

\textbf{Comprehensive validation} (\texttt{comprehensive\_validation.json}):

\textbf{Hardware validation \texttt{validation\_passed: true}} (OBU anchor); extrapolation block reports \textbf{stable} energy-per-second across durations tested.

\subsection*{Newest sweep: \texttt{parameter\_sweep.json}}

\begin{itemize}
\item \textbf{108} grid combinations (EA thresholds + PHY-style knobs).
\item \textbf{Best} observed \textbf{energy saving vs RSSI (\%)} in that grid: \textbf{$\sim$6.90\%}, at \textbf{\texttt{bs\_coverage\_radius}: 400}, \textbf{\texttt{tx\_power\_max\_watts}: 0.75}, \textbf{\texttt{sinr\_handoff\_trigger\_db}: 1.0}, \textbf{\texttt{max\_acceptable\_epb}: 1.5$\times$10$^{-7}$}, \textbf{\texttt{hysteresis}: 0.15} (plus availability/outage/handoffs per row in the file).
\end{itemize}

\subsection*{Short answer}

The \textbf{latest simulated artifact} is \textbf{\texttt{parameter\_sweep.json}} (connectivity/EA hyperparameter sweep). The \textbf{latest full baseline snapshot} in the repo is \textbf{\texttt{simulation\_results.json} for seed 1011}, showing \textbf{$\sim$3\% EPB} and \textbf{$\sim$1.36\% total energy/CO\textsubscript{2}} improvement vs RSSI, with \textbf{$\sim$11\%} fewer handoffs; \textbf{scaling} and \textbf{sensitivity} runs show \textbf{roughly 3--5\%} EPB savings depending on scenario, and \textbf{scenario\_diversity} / \textbf{weather\_sweep} add condition-specific numbers (including \textbf{very large} savings in \textbf{heavy rain} in the diversity setup).
\end{assistmsg}

% ========== TURN 3 (this request) ==========
\begin{usermsg}
Now check our this whole chat.

Then make this whole chat into a detailed good looking latex.

FUll chat but in a detailed overview type latex(cool and good looking).

But NO TEXT OR ANY WORDS ALTERATION.

KEEP EVERYTHING WHAT IS IN THIS CHAT BUT IN A OVERVIEW OF THIS PROJECT ALIKE.
\end{usermsg}

\end{document}
